\documentclass{ctexart}
% claim-forward 四类仅 LLM 处置的中文合成 fixture（latex-thesis-zh）。
% Case A–D 为四类处置的正例（预设反驳壳、让步后否定、断言与让步壳共现、空泛谦辞）；
% Case E–H 为应保持不动的边界。供 evals/evals.json id 61 使用。
% 四类处置没有 CF-* 码，脚本测试不读本文件。
\begin{document}

\chapter{绪论}
本文研究预约提醒流程对社区慢病随访的影响。

% Case A：预设反驳壳——回应了数据与所引文献都未提出的异议，回应本身没有证据、条件或范围。
本文提出的预约提醒流程将平均候诊时间由42分钟缩短至31分钟（表~\ref{tab:wait}）。
这并不意味着该流程能够解决社区卫生服务中的所有问题。

\chapter{数据与方法}
表~\ref{tab:interview}汇总了16名社区医生的访谈，其中10名把“居民联系方式变更后未及时登记”列为失访的主要原因。
表~\ref{tab:age}显示，80岁以上居民的按期随访比例为68.0\%，低于全体居民的80.0\%。

% Case F：边界——定义性否定，用于区分两个口径。
本文所称“按期完成随访”指在预约日前后7日内完成面访或电话随访，不包括仅发送短信提醒的情形。

\chapter{实验结果}
% Case B：让步后否定——分母发生变化，主句不给分析就否定其影响。
虽然随访登记人数由250增至320，按期完成随访的比例由84.0\%降至80.0\%，但这并不影响新提醒流程提高了随访依从性的判断。

% Case C：断言与让步壳共现——质性证据上的过强断言加模糊让步。
访谈结果充分证明联系方式变更是失访的决定性因素。
当然，不能简单地认为所有社区都如此，但总体而言该因素在一定程度上起到了关键作用。

% Case E：边界——带原因的真实证据边界。
访谈记录的是社区医生自述的失访原因。
受访者均来自同一城区，结果不说明各原因在其他城区的实际占比。

\chapter{讨论}
% Case G：边界——对带引用的前人工作的让步，并写明差异来源。
虽然已有研究\cite{followup2020}报告了更高的随访完成率，但其统计口径只包含能够到院复诊的居民。

% Case D：空泛谦辞——限制句没有对象、条件或量值。
本文提出的流程仍存在一定不足，仍有提升空间。

\chapter{结论与展望}
本文提出的预约提醒流程将平均候诊时间由42分钟缩短至31分钟。

% Case H：边界——写出具体局限的承接句，其后接展望。
本文流程在80岁以上居民中的按期随访比例（68.0\%）仍低于全体居民（80.0\%），且访谈只覆盖一个城区。
后续工作将把家属联系人纳入提醒对象，并在第二个城区重复本文的访谈。

\end{document}
